% Pillar-4 (length bias) standalone abstract
\begin{abstract}
A recurring worry about GRPO is length bias: because the group-relative update
divides by response length, the algorithm may be nudged toward longer completions
that stumble into reward rather than toward better reasoning. The length-debiased
variant Dr.\,GRPO removes this per-response normalization. We study length bias and
held-out generalization as one pillar of \ours{}, a benchmark of 70+ RL runs across
seven libraries and five model families ($0.6$B--$\sim$671B) on GSM8K, HumanEval,
and tool-use tasks. Our headline result is a carefully scoped negative one: in the
regime we measure---short-horizon GSM8K chain-of-thought with a 200-token
generation cap---the mean Qwen3-8B GRPO
gain over the pre-RL checkpoint on held-out data is small and not statistically
significant, GRPO and Dr.\,GRPO produce held-out gains that are indistinguishable
at our seed budget, and neither inflates length within this capped regime (pooled
mean completion length drifts from $194.4$ to $183.5$ tokens). We stress that the
$200$-token cap \emph{bounds} this claim: it cannot rule out length inflation in
longer-horizon or uncapped regimes, only establish its absence where length is
already controlled. The verbosity-trap signature is absent \emph{in this regime}. We argue this is the
\emph{expected} outcome when length is already controlled, not evidence against
Dr.\,GRPO, and specify the length-confounded regime and causal mediation tests that
would actually reveal a difference. We release code, logs, and checkpoints.
\end{abstract}
\paragraph{August 2026 evidence boundary.} This bounded null is not evidence
of a general deployment or capacity claim outside its named length-capped setup.
